\chapter{Lowering and Execution}
\label{ch:pebble-lowering}

An interpreter can implement a phi directly because it remembers the predecessor block.
A simpler register machine should not need to know what a phi is.
Lowering turns path-dependent selection into ordinary moves executed on the selected edge, then translates block targets into instruction addresses.

\section{A target with explicit storage}
\label{sec:pebble-vm}
The VM has input, constant, add, multiply, compare, move, jump, branch, and return instructions.
Registers 1 through 128 correspond to IR value IDs.
Registers 129 through 256 are scratch locations used to sequence edge copies.
There is no physical register allocation: this is a virtual-register target with enough bounded storage for every value.
The target is useful for validating lowering before instruction constraints and spills complicate the argument.

The VM decoder is not a serialized-bytecode loader.
\code{pir\_vm\_run} executes an initialized \code{PirVm} normally produced by verified lowering.
It checks opcode, register, immediate, and target bounds while executing, but it does not independently prove that arbitrary hand-written VM code defines each register before use.
Those static properties are guaranteed by the lowering contract for its output.

\section{Copy on the edge, not on both paths}
\label{sec:pebble-edge-lowering}
For each edge $P\to B$, collect the incoming operand of every phi in $B$ labeled by $P$.
The required assignment is a parallel bundle from those sources to the phi result registers.
Algorithm~\ref{alg:pebble-edge-copies} sequences the bundle without cycle detection by using one scratch slot per destination.
It spends extra moves to make the proof simple.

\begin{algorithm}[htbp]
\caption{Lowering a phi bundle into two phases of moves.}
\label{alg:pebble-edge-copies}
\begin{pseudocode}
\Require verified edge $P\to B$, distinct scratch register $s_v$ for each phi $v$ in $B$
\Ensure sequential moves implement the edge's parallel assignment
\ForAll{phis $v$ in $B$}
\State emit $s_v\gets$ the source register selected by predecessor $P$
\EndFor
\ForAll{phis $v$ in $B$}
\State emit $v\gets s_v$
\EndFor
\State emit a jump to $B$'s first ordinary target instruction
\end{pseudocode}
\end{algorithm}

All sources are read before any destination is overwritten, so the resulting bundle preserves swaps and longer cycles.
For the swap fixture, the back edge first saves old \code{a} and old \code{b}, then publishes both new values.
Two loop iterations return the initial pair rather than collapsing it to one repeated value.
The test exercises the IR interpreter and the lowered VM against the same direct swap expectation.

For a conditional branch, lowering emits two edge-specific stubs.
The VM branch selects one stub, which performs only that edge's copies and jumps to its destination.
Thus even a critical edge has its own execution point without modifying the source IR graph.
Figure~\ref{fig:pebble-edge-stubs} shows the arrangement.

\begin{figure}[htbp]
\centering
\begin{tikzpicture}[node distance=11mm and 10mm]
\node[basic block] (b) {VM branch};
\node[basic block,below left=of b] (t) {True-edge copies};
\node[basic block,below right=of b] (f) {False-edge copies};
\node[basic block,below=of t] (bt) {True target body};
\node[basic block,below=of f] (bf) {False target body};
\draw[cfg edge] (b)--(t); \draw[cfg edge] (b)--(f);
\draw[cfg edge] (t)--(bt); \draw[cfg edge] (f)--(bf);
\end{tikzpicture}
\caption{Target-side edge stubs execute only the copy bundle selected by the branch.}
\label{fig:pebble-edge-stubs}
\end{figure}

\section{Addresses after layout}
\label{sec:pebble-fixup}
Ordinary instructions are emitted block by block.
Each block records its first target instruction index, before its ordinary instructions or terminator are emitted.
Forward jumps temporarily hold a block number.
After all blocks are laid out, a fixup pass replaces every such number with the recorded instruction index.
Conditional stub addresses are known during emission and already use instruction indices.
The distinction is internal to the builder; no partly fixed-up VM is published.

For the optimized running expression there are five value-producing instructions and one return, hence six dispatches.
The target returns 128 for input 9 at budget six and exhausts at budget five.
This test checks the target's fuel boundary separately from the IR interpreter's one-block boundary.

\section{What comes after this target}
\label{sec:pebble-target-cost}
For $N$ allocated IR slots, $B$ blocks, and $P$ live phis, the current scans cost $O(BN)$.
Each phi has two incoming edges and produces two moves per edge, so edge copying emits $4P$ moves.
Ordinary instructions contribute at most $N-P$, and terminator scaffolding contributes at most $3B$ instructions.
The loose total bound is therefore $N+3P+3B$, before any target cleanup.
Capacity checks still guard every emission so changes to these assumptions cannot silently corrupt storage.

A physical target would add register classes, allocation, spills, calling conventions, and machine encodings.
It should reuse the verified edge-copy semantics rather than inventing a second interpretation of phi selection.
Linear scan is a reasonable next allocation experiment, but its interval and instruction-constraint assumptions need their own contract~\cite{poletto-1999-linear-scan}.

\section{Exercises}
\label{sec:pebble-fifteen-exercises}
\begin{enumerate}
\item \label{ex:pebble-copy-bundle} Sequence the parallel bundle $a\gets b$, $b\gets a$, $c\gets a$ using distinct scratch locations. Trace it from $(a,b,c)=(4,9,0)$.
\item \label{ex:pebble-fixup} A forward jump names block 3, which is eventually laid out at instruction 27. Explain what must be changed at fixup and why block 3 is not itself a valid final address by convention.
\end{enumerate}

\section{Further reading}
\label{sec:pebble-fifteen-reading}
Appel provides a complete backend context for the transition from abstract values to machine storage~\cite{appel-1998-modern}.
Poletto and Sarkar describe linear scan's allocation strategy and its compilation-cost motivation~\cite{poletto-1999-linear-scan}.
